\section{Related Work}
\label{sec:related}

\paragraph{Narrative Framing in Online Discourse.}

Framing, the process of selecting and highlighting parts of reality to shape interpretation~\citep{Entman1993}, is a central mechanism of persuasion. Narratives make frames concrete and memorable. \citet{Green2000} show that being absorbed in a story reduces critical evaluation and promotes story-consistent belief change, a finding supported by broader evidence on narrative persuasion~\citep{BraddockDillard2016}. In online social media, such strategies have become scalable tools for opinion manipulation~\citep{Tucker2018,Wardle2017}. Different narratives draw on different mechanisms: victim stories invite empathetic identification~\citep{Green2000}, conspiracy theories offer coherent causal explanations under uncertainty~\citep{Douglas2017}, fear appeals heighten threat perception~\citep{Witte1992}, and authority-based claims exploit credibility shortcuts~\citep{Cialdini2001}. Recent NLP work has put these categories into practice through propaganda technique detection~\citep{DaSanMartino2019}, multi-dimensional propaganda analysis~\citep{liu-etal-2025-propainsight}, narrative framing models for political discourse~\citep{otmakhova-frermann-2025-narrative}, and large-scale resources for hyperpartisan and conspiratorial narratives across languages~\citep{maggini-etal-2026-partisanlens}. Most of this work, however, studies manipulation retrospectively as text analysis. Less is known about how narrative strategies differ once they enter repeated, multi-party discussion.

\paragraph{Computational Approaches to Online Manipulation.}

Computational studies of online manipulation have mainly followed detection and modelling paradigms. Detection work identifies propaganda techniques~\citep{DaSanMartino2019}, mines disinformation narratives~\citep{sosnowski-etal-2025-dinam}, integrates inoculation theory into LLM-based disinformation detection~\citep{modzelewski-etal-2026-malicious}, and profiles bot-driven or coordinated campaigns~\citep{Ferrara2016,Woolley2018}. Modelling work, especially agent-based opinion dynamics~\citep{Deffuant2000,Flache2017}, captures convergence and polarisation under social influence, but usually treats opinions as scalar states and leaves out the language through which influence occurs. Meanwhile, LLMs have been shown to be persuasive communicators. LLM-generated messages can shift policy attitudes as effectively as human-written arguments~\citep{Bai2025}, GPT-4 can outperform human debaters when personalised~\citep{Salvi2025}, model persuasiveness appears to scale with capability~\citep{pauli-etal-2025-measuring}, and demographic-aware human--LLM interaction can drive echo-chamber and polarisation patterns~\citep{borah-etal-2026-persuasion}. These studies establish the risk of LLM-driven persuasion, but mostly examine single-message or two-party interaction. We instead study manipulation inside a networked discussion where messages accumulate, circulate, and reshape later responses.

\paragraph{LLM-Based Social Simulation.}

LLM-based social simulation provides a controlled way to study collective behaviour. \citet{Park2023} introduced memory, reflection, and planning for believable generative agents. Recent studies use such agents to examine debate biases~\citep{taubenfeld-etal-2024-systematic}, group conformity~\citep{choi-etal-2025-empirical}, echo chambers~\citep{wang-etal-2025-decoding}, and stance--topic co-evolution~\citep{zhang-etal-2025-spark}. Others study content spread and moderation~\citep{liu-etal-2025-mosaic} and news distortion~\citep{liu-etal-2025-stepwise}. Recent critical work also stresses that the validity of such simulations depends on empirically testing how realistically generative agents reproduce online communication~\citep{munker-etal-2026-dont}, and on keeping open-ended free-text outputs rather than reducing opinion to multiple-choice probes~\citep{ma-etal-2026-open}. These works show that LLM agents can support empirical inquiry into social dynamics, including misinformation-related settings. What remains less explored is the comparative effect of distinct narrative attack strategies, their interaction with population traits, and their response to interventions. NarrateSim is a social-simulation framework built for this question, serving as an empirical testbed for narrative attacks across topics, populations, network positions, and intervention conditions.
